\section{QR algorithm}

\subsection{Task}

Implement the QR decomposition of a matrix. Based on it, implement the QR algorithm for the full
eigenvalue problem of an arbitrary matrix taking the matrix and the precision as input, and find
the eigenvalues of the matrix.

\[
A = \begin{pmatrix}
6 & -3 & 7 \\
9 & 1 & -6 \\
3 & -5 & 5
\end{pmatrix}
\]

\subsection{Method}

The QR decomposition $A = QR$ with an orthogonal $Q$ and an upper triangular $R$ is built with
Householder reflections. At step $k = 1, \ldots, n - 1$ the vector $v$ is formed from the part of
column $k$ on and below the diagonal:

\[
v = a_k + \operatorname{sign}(a_{kk}) \|a_k\|_2 e_k, \qquad
H_k = E - \frac{2 v v^T}{v^T v}.
\]

The reflection $H_k$ zeroes the elements below the diagonal in column $k$, and the choice of the
sign avoids cancellation. Then $R = H_{n-1} \cdots H_1 A$ and $Q = H_1 \cdots H_{n-1}$. The
reflections are applied as rank-one updates without forming $H_k$.

The QR algorithm builds the sequence $A^{(0)} = A$, $A^{(k)} = Q^{(k)} R^{(k)}$,
$A^{(k+1)} = R^{(k)} Q^{(k)} = Q^{(k)T} A^{(k)} Q^{(k)}$ of similar matrices. It converges to a
block upper triangular form: $1 \times 1$ blocks hold real eigenvalues and $2 \times 2$ blocks hold
pairs of complex conjugate eigenvalues. The diagonal is scanned column by column:

\begin{itemize}
    \item if $\bigl(\sum_{i>j} a_{ij}^2\bigr)^{1/2} \le \varepsilon$, then $a_{jj}$ is a real
    eigenvalue;
    \item otherwise columns $j, j + 1$ form a block, whose eigenvalues are the roots of
    $(a_{jj} - \lambda)(a_{j+1,j+1} - \lambda) = a_{j,j+1} a_{j+1,j}$. The block is accepted when
    the elements below it are smaller than $\varepsilon$ and its eigenvalues differ from those of
    the previous iteration by at most $\varepsilon$.
\end{itemize}

The iterations stop when every diagonal position is covered by an accepted block. For complex
eigenvalues the sub-diagonal element of the block does not vanish, so the stability of the block
eigenvalues is the only possible criterion.

\subsection{Results}

The program reads $n$, the matrix $A$ and the precision $\varepsilon$ and prints the QR
decomposition of $A$, the number of iterations and the eigenvalues.

\needspace{10\baselineskip}
\lstinputlisting[style=console, caption={tests/data/qr/01.in}]{../tests/data/qr/01.in}
\lstinputlisting[style=console, caption={tests/data/qr/01.out}]{../tests/data/qr/01.out}

The matrix has one real eigenvalue and a pair of complex conjugate eigenvalues. The unit tests
check $QR = A$ and $Q^T Q = E$, compare the eigenvalues of symmetric matrices with the rotation
method and check that the sum and the product of the eigenvalues are equal to the trace and the
determinant of random matrices up to $6 \times 6$.

\subsection{Analysis}

\begin{table}[!ht]
\centering
\begin{tabular}{|r|r|r|r|r|r|}
\hline
\rowcolor{lightgray}
$\varepsilon$ & $10^{-1}$ & $10^{-2}$ & $10^{-3}$ & $10^{-4}$ & $10^{-8}$ \\
\hline
Iterations & 7 & 9 & 12 & 15 & 26 \\
\hline
\end{tabular}
\caption{Number of QR iterations depending on the precision}
\end{table}

The elements below the $2 \times 2$ block decrease as
$\bigl(|\lambda_3| / |\lambda_{1,2}|\bigr)^k = (3.83 / 8.81)^k \approx 0.435^k$, which gives
$1 / \lg(1/0.435) \approx 2.8$ iterations per decimal digit, in agreement with the table.

\subsection{Source code}

\lstinputlisting[caption={include/qr.hpp}]{../include/qr.hpp}
\lstinputlisting[caption={src/qr.cpp}]{../src/qr.cpp}

\pagebreak
